% ED26-BM-RAT2021-3-2 C01--C05: exact O'Brien2021 cognitive account.
\chapter{Natural logic: representations, schemas and pragmatic control}
\label{sec:p19-rationality-3-2}

O'Brien's \emph{Natural Logic} is a cognitive account of how people record
propositions and construct inferences, not a new name for classical
natural deduction. His selected case is the Braine--O'Brien mental-logic
theory. Its explanatory package has three parts: a representation with
predicate--argument structure, inference schemas together with a program
controlling their use, and pragmatic principles integrating inference
with knowledge and practical aims. Merely listing valid formal rules
supplies none of the other two parts
\cite[pp.~215--217]{OBrien2021NaturalLogic}.

The source argues for that package; it does not prove a completeness
theorem for human cognition. Its historical narrative contrasts an older
identification of logic with thought, including Boole, with later
anti-psychologism and psychology's separation from philosophy. O'Brien
regards that separation as a change of intellectual attitude rather than
experimental evidence that reasoning lacks logic. This is his historical
and methodological argument, not a controlled comparison establishing
that one contemporary architecture uniquely describes every reasoner.

\section{What the proposed architecture must explain}
\label{sec:p19-r32-architecture}

The argument begins from declarative memory: recording what an entity
does or what properties it has requires a format that preserves relations
between entities and predicates. Propositional attitudes additionally
distinguish assertion, denial, supposition, questioning, belief and doubt.
O'Brien proposes a logical language of thought for this structure and
inference procedures that operate on it. A format alone cannot explain
which new conclusions are produced, and a stock of schemas alone cannot
explain when a reasoner applies them. Hence the program and the integration
with practical information are indispensable to the proposal
\cite[pp.~215--217]{OBrien2021NaturalLogic}.

The source illustrates an evolutionary motivation with a gatherer who
knows that food is available near a river or near a canyon. Information
that the river location is unproductive leads toward the canyon. The
example motivates elimination of an unavailable alternative; it is not
an observed ancestral experiment or proof that natural selection installed
these exact schemas. Likewise, the fact that a representation needs
relational structure does not alone establish that the Braine--O'Brien
schema inventory and control program are its only possible implementation.
The source's stronger necessity claim must be evaluated by comparing
the actual alternative mechanisms and their predictions.

O'Brien explicitly separates this project from formal system construction.
A cognitive theory may omit some textbook-valid inferences and explain
why people fail or decline to make them. Completeness was not a proposed
evolutionary goal; soundness matters, but interpretation and admissible
background information must first be fixed. The finite formal benchmarks
below test particular inferences under declared readings. They do not
turn psychological procedures into a sound-and-complete calculus.

\section{Representation and inferential-role meaning}
\label{sec:p19-r32-representation}

The source uses language-of-thought expressions rather than identifying
surface English words with truth-table connectives. Its \emph{if} is not
defined as the material conditional. Its illustrative material-implication
paradox combines not having a gun with a conditional admission of guilt:
a false antecedent makes a material implication true, but this does not
make that utterance a plausible ordinary conditional. Conversely a threat
worded with \emph{and}, such as reaching into a cookie jar and receiving
a punishment, can convey a conditional rather than two asserted facts.
Thus input must be interpreted into the internal representation, and a
conclusion must be expressed back in ordinary language. Surface matching
of the word \emph{and} is not the interpretation algorithm
\cite[pp.~216--217]{OBrien2021NaturalLogic}.

Meanings in the logical vocabulary are explained by the inferences that
the terms sanction. The chapter calls this conceptual-role or
inferential-role semantics and explicitly limits its own claim to the
natural-logic vocabulary, rather than providing a complete semantic theory
for every lexical item. Consider the source's zoo example: all children
who visited the zoo were in the third grade versus some such children
were. Learning that Álvaro visited licenses the grade conclusion from
the universal but not from the existential. Understanding that contrast
is part of the proposed grasp of \emph{all} and \emph{some}.

\begin{hbtheorem}{thm:p19-r32-quantifier-benchmarks}{proved=annex:p19.r32.quantifiers}{Pointwise quantified benchmarks and their limits}
Under the ordinary equality-free first-order reading,
\(\forall x(Zx\to Gx),Za\) entails \(Ga\), whereas
\(\exists x(Zx\land Gx),Za\) need not.
For predicates \(X,A,P,Q\), the three premises
\[
 \forall x(Xx\to(Px\lor Qx)),\quad
 \forall x(Ax\to Xx),\quad \forall x(Ax\to\neg Qx)
\]
entail \(\forall x(Ax\to Px)\).
Replacing the third premise by \(\neg\forall x(Ax\to Qx)\)
does not preserve that conclusion. These are specified formal
benchmarks, not the definition of the source's cognitive notation.
\end{hbtheorem}

Quantifiers in the source representation are content-bound. Its example
\(S_1[\mathrm{All}\ X]\ \mathrm{OR}\ S_2[\mathrm{PRO\!-!All}\ X]\)
uses the pronoun-like PRO annotation to carry the first quantifier's scope
into the second clause. The chapter contrasts this with a quantifier
placed outside a textbook formula. It reports Ioup's survey of fourteen
languages as linguistic motivation; it supplies neither the language
selection protocol nor the evidence needed to infer a universal claim
about every language or about cognitive implementation. That dependent
source extent remains distinct from an experiment on fourteen people
\cite[p.~217]{OBrien2021NaturalLogic}.

The corresponding illustrated inference combines that quantified
alternative with a selected group \(\alpha\) inside \(X\) and a
negative \(S_2\) judgment about \(\alpha\), yielding \(S_1\) for that
group. Uirá and Werá, boys of a village, are not fishing by the river, so
the stated alternatives locate them hunting in the forest. The theorem's
pointwise reading makes its validity check explicit. It must not identify
``not all members satisfy \(S_2\)'' with ``all members fail \(S_2\)''.

Collective and distributive arguments are another material distinction.
The source contrasts boys jointly carrying one box with a distribution
of box-carrying over boys, using a group argument versus an All-marked
argument. Its prosecutors--jurors example makes the prosecutors a
collective argument while their persuasion reaches individual jurors.
A pointwise \(\forall\exists\) translation by itself does not require
distinct boxes for different boys; that stronger reading needs an
explicit distinct-witness condition. Nor does collective persuasion
license the assertion that every prosecutor individually persuaded
every juror. A surface sentence therefore does not automatically fix
one quantified structure \cite[pp.~217--218]{OBrien2021NaturalLogic}.

\section{Schemas are not the whole reasoning procedure}
\label{sec:p19-r32-program}

The chapter describes a sentential inventory advertised as fourteen
schemas, referring to Braine and O'Brien's 1998 Tables~6.1--6.2 for the
full presentation. Seven are described as core schemas, applied whenever
their conditions are met. The illustrated examples are disjunction
elimination from an alternative and a negated disjunct, and modus ponens
from a conditional and its antecedent. Feeder schemas build a conjunction
from available components and select a component from a conjunction,
but only when their output feeds another inference. This restriction
explains why conjunction construction and projection can support an
inference without being gratuitously reported as new conclusions or
repeated in loops. Two further schemas concern incompatibility and two
supposition: conditional proof and a contradiction-based rejection of a
supposition \cite[p.~218]{OBrien2021NaturalLogic}.

These are source-described families and examples, not a fully reproduced
fourteen-row calculus. The paragraph counts seven core, two feeder
families and four further schemas; that is thirteen families, while
projection has two orientations. Whether the source tables count those
orientations separately must be checked in those exact tables. This
entry does not invent a fourteenth row or silently combine the 1998
inventory with Rips's different theory. In particular, the chapter says
Braine--O'Brien excludes disjunction introduction from one disjunct on
empirical grounds, whereas Rips includes it. Validity in classical logic
does not settle a cognitive inventory or conclusion-report policy.

\begin{hbtheorem}{thm:p19-r32-sentential-benchmarks}{proved=annex:p19.r32.sentential}{The source's chained example and a conditional-evaluation contrast}
Under the Boolean benchmark, the source's premises
\(N\lor P,\neg N,P\to H,H\to Z,\neg(Z\land S)\)
entail, in succession, \(P,H,Z,\neg S\). The same consequences follow
when the premises are listed in reverse order. Also \(p\) entails
\(p\lor q\), whether or not a cognitive theory includes that
conclusion-producing schema.
By contrast, derivability of \(\neg q\) under an added \(p\)
supposition does not establish that a material \(p\to q\) was
false under the original information: a valuation with both \(p,q\)
false is a counterexample to that proposed test.
\end{hbtheorem}

The direct-reasoning routine, DRR, is proposed to be automatic and
universally available; further strategies are described as widely
available among college students, not as part of that universal
routine. Conditional evaluation supposes \(p\) and reasons toward
\(q\): obtaining \(q\) supports the conditional, obtaining
\(\mathrm{not}\ q\) rejects it. If neither is obtained, the described
routine alone has not settled the conditional. This evaluation procedure
is expressly different from a material-implication truth test. The
conditional-proof and modus-ponens procedures are intended to work
together: a consequent previously dependent on a supposition can be
asserted without that marking when the antecedent is subsequently known.
This is a claim about cognitive control and inferential meaning, not a
new classical truth function \cite[pp.~218--219]{OBrien2021NaturalLogic}.

Suppositional reasoning also needs an admissibility policy. Information
used under a counterfactual supposition must still be applicable under
it. The source illustrates the problem with Philadelphia lying north
of New York: changing the geography leaves the consequences of moving
either city underspecified. It proposes negotiation over which
background propositions to exclude, rather than requiring a closest-world
selection. This identifies an unresolved interpretation/control choice;
it is not an algorithm proving that closest-world theories are impossible.
The chapter's short DRR description likewise supplies no conflict policy
for deriving both \(q\) and its negation from inconsistent retained
assumptions, and is not a complete implementation of every referenced
strategy or predicate-level schema.

\section{Pragmatic integration without changing formal standards}
\label{sec:p19-r32-pragmatics}

An offer to pay for painting a shed can invite a biconditional reading,
including no payment without painting; the source distinguishes this
invitation from the logical content of a one-way conditional. Context can
also favor inclusive or exclusive alternatives, without the natural-logic
inventory itself fixing every utterance to one of them. A diner script
may make choosing both tea and coffee unlikely and feed that information
into an otherwise logical inference. The source proposes that people
often do not tag the varied origins of their inferences in ordinary
reasoning \cite[p.~219]{OBrien2021NaturalLogic}.

Such integration explains why a response beyond a restricted formal
premise set need not be irrational under the participant's richer
interpretation. It does not make that response follow from the restricted
premise set, or make every script reliable. The permission to consider
practical information, the truth of that information, and the validity
of an inference given it remain different questions. The existing
conditional suppression and interpretation treatments make this
distinction consequential for actual tasks
(\cref{sec:p19-conditional-reasoning-suppression,sec:p19-conditional-interpretation-experiments}).

\section{What the empirical claims actually supply}
\label{sec:p19-r32-evidence}

The source predicts that DRR-accessible inferences will usually be made,
that textbook-valid but unavailable schemas and extra-strategy problems
will be harder, and that schema availability will explain difficulty
and order beyond problem length. It reports additive schema weights
estimated from difficulty ratings and tested against new problems and
new participants: 66\% variance for sentential problems and 69\% for
predicate problems, becoming 55\% and 56\% with problem length partialed
out. These are chapter-reported fit quantities attributed to Braine,
Reiser and Rumain 1984 and Yang, Braine and O'Brien 1998. They are not
percentages of people using a mental logic. Their exact samples,
regression units, coefficients and uncertainty are not supplied by the
selected chapter or recovered here, so no independent replication or
unique-process inference is claimed \cite[pp.~219--220]{OBrien2021NaturalLogic}.

The source also attributes inference-order evidence to O'Brien, Braine
and Yang 1994 and Braine and colleagues 1995. Its five-premise chain
produces the same proposed conclusion order in normal and reversed
presentation orders. The formal theorem above verifies consequences,
not the empirical frequency or timing of that sequence. The chapter
does not provide participant counts, order-assignment details or
uncertainty for that claim. Its claim that core schemas occur in young
children depends on Braine and Rumain 1983 and selected chapters of the
1998 volume; those child protocols and samples have not been reconstructed
from the 2021 discussion.

\subsection{A recovered primary text-comprehension study}

O'Brien, Roazzi, Dias and Soskova's Portuguese-language 2007 paper provides
three distinct UFPE student experiments, not a population-wide study of
spontaneous reasoning. Experiment~1 recruited twenty students, five male
and fifteen female, ages 19--24, mean 20.4, SD 1.4. They judged paired
valid/invalid endings for sixteen short Portuguese stories, using eight
predicate-schema types with parallel content versions. Presentation and
ending orders were controlled, and testing occurred in small groups.
The reported accuracy was 93\% before exclusion. Three participants with
four, five and seven errors were removed as apparent outliers, leaving
seventeen and reported accuracy 97\%. Both denominators matter: the
exclusion is not evidence that the recruited sample originally had 97\%
accuracy \cite[pp.~121--122]{OBrienRoazziDiasSoskova2007Text}.

Experiment~2 used twenty-nine psychology undergraduates, three male and
twenty-six female, mean age 22.1, SD 3.4. After sixteen stories without
explicit final conclusions, participants judged whether a model-predicted
inference, a paraphrase, and a contrast inference had been presented,
were uncertain, or had to be inferred. The order was randomized, and
participants were told not to turn back to the story. The paper controlled
word lengths, negations, textual similarity, premise positions and
separation. Coding the three responses 1, .5 and 0 gave means .65 for
model-predicted items, .81 for paraphrases and .27 for contrasts. Recoding
uncertainty as 0 gave .63, .78 and .14. These are item-recognition scores,
not proportions of participants whose brains implement a particular
rule. The Friedman statistic is reported as \(\chi^2(2)=42.11\);
its printed bound \(p<.000\) is mathematically impossible and is not
repeated here as a valid probability claim
\cite[pp.~122--124]{OBrienRoazziDiasSoskova2007Text}.

Experiment~3 recruited thirty-eight psychology undergraduates, nine male
and twenty-nine female, mean age 21.0, SD 2.8. Nine stories of five to ten
sentences required two or three predicate schemas in succession; each
was followed by an ending choice and then the recognition task. Two
story orders and recognition-item orders were randomized. Mean correct
ending choice was .92, SD .10. Recognition responses for stories with
an incorrect ending choice were excluded, so that analysis is conditional
on successful initial performance. The dichotomized recognition means
were .58 for model-predicted items, .77 for paraphrases and .17 for
contrasts. The reported Friedman statistic was \(\chi^2(2)=65.51\),
\(p<.001\), with paired Wilcoxon comparisons favoring the two successive
contrasts. The paper supplies no confidence intervals for these means
\cite[pp.~124--126]{OBrienRoazziDiasSoskova2007Text}.

The results establish that these participants often selected the intended
endings and treated certain unpresented conclusions as presented
information. The authors interpret this as easy, routine inference during
reading. Recognition/source-memory judgments do not by themselves
identify a unique internal rule sequence or measure zero cognitive effort.
Experiment~3's recognition task followed an explicit ending-choice decision,
so it is not independent evidence of inference before any prompting.
Comprehension strategies, reconstruction at judgment, lexical cues and
alternative model procedures remain relevant competitors. The comparison
of single- and multiple-schema experiments was not randomized assignment
of one representative population to formal predicate and sentential
calculi. It therefore does not establish equal difficulty for all tasks
at those levels.

The finite outstanding evidence set also includes Lea and colleagues
1990 and Lea 1995 for sentential reading, Ioup 1975 for language structure,
the source's named child-schema extents, and the original order/weight
studies. The 2007 study corroborates a particular predicate-reading task;
it does not substitute for those distinct protocols. Their unrecovered
sample/design/uncertainty information remains an explicit source gap,
not an invitation to invent missing observations.

\section{Objections to competing accounts and their proper reach}
\label{sec:p19-r32-comparisons}

O'Brien argues that content-specific permission/obligation schemas and
social-contract accounts explain too narrow a range of tasks to replace
a general inference architecture. He asks how many content domains would
be needed, why ordinary conditionals convey these rules, and whether
children's early conditional competence fits an account in which general
rules arise only after many acquired pragmatic schemas. His methodological
concern supplies no metric or proved lower bound on the number of modules.
The further
contrast is between falsifying an indicative conditional, a metalevel
truth-testing task, and identifying a violation of a regulation. Better
performance with permission or social-exchange wording does not isolate
one representation if instructions, interpretation and practical aims
also change. The chapter's comparative claims about Cheng--Holyoak and
Cosmides require their exact studies, not a universal ranking inferred
from the labels \cite[pp.~220--221]{OBrien2021NaturalLogic}.

For Bayesian accounts the source argues that neutral, unfamiliar content
does not supply the prior beliefs on which conditional acceptance is said
to depend, and asks for predictions across conjunction, negation,
disjunction and conditional conclusions. This is a demand for a specified
model and wider task coverage. Lack of prior \emph{content knowledge},
however, is not a mathematical impossibility of assigning a probability
model. Nor does success on a wider task set establish that each rival
can only be extended by adopting these particular schemas. The exact
PHM/PRM task and fitting boundaries in
\cref{sec:p19-rationality-4-5} must not be replaced by the source's broad
comparative rhetoric.

Evans and Over's conditional-probability proposal receives a more specific
integration argument: supposing the antecedent and considering consequent
versus negated consequent resembles conditional proof, but probability
assessment adds information not supplied by a deductive schema alone.
O'Brien illustrates this with a red box containing eighty beads and a
blue box containing twenty, suggesting .80 when the known bead comes
from one of them. The example is an explanatory illustration, not an
experiment with a recruited sample.

\begin{hblemma}{thm:p19-r32-bead-sampling}{proved=annex:p19.r32.beads}{Counts alone do not fix the bead probability}
If each of the hundred beads is equally likely to be selected,
the probability of red origin is \(80/100=.8\). If instead either
box is selected with probability \(1/2\) and then a bead uniformly
within that box, red origin has probability \(1/2\).
Thus the two box sizes alone do not determine the illustrated probability;
a selection law is needed. This is a stipulated calculation, not an
estimate of people's judgments.
\end{hblemma}

The source's mental-model objection concerns representations with no
retained variables or quantifiers. It discusses a conditional model with
an exhaustivity marker and an ellipsis, the combination with its
antecedent to obtain the consequent, and quantified models constructed
from parsed premises by recursive loops. Its arithmetic illustration
starts from every selected \(x\) being a sum of some \(y,z\), but
ends in finite lists of exemplars. O'Brien asks how one could recover
the original quantified proposition, even for reiteration, if the
source syntax and scope have disappeared. Exhaustivity markers and
retained premodel representations are material: a finite unannotated
sample is not the same as a model carrying explicit domain exhaustion.

The complaint is a representational challenge, not a proof that all model
theories are unable to reason. The current 2021 mental-model treatment in
\cref{sec:p19-rationality-2-3} includes meaning, symbolic operations and
counterexample search; it is not identical to the source's selected
1991/1993 descriptions. O'Brien's own conclusion allows counterexample
search and specialized violation detection as valuable additions. He
opposes making those special procedures the whole architecture, rather
than denying their utility \cite[pp.~222--223]{OBrien2021NaturalLogic}.

\begin{hbtheorem}{thm:p19-r32-generalization}{proved=annex:p19.r32.generalization}{Finite exemplars and the decimal example}
Finite positive records of a relation, without a declared exhaustive
domain or retained universal premise, do not in general determine a
universal \(\forall\exists\) assertion. This qualification does not
remove an explicitly supplied exhaustivity marker.
Separately, every natural number whose decimal representation ends in
zero is divisible by five, including zero itself. The latter has a
general arithmetic argument, not merely verification of a few examples.
\end{hbtheorem}

\section{Consultation and evidence boundaries}

Use this account to compare a proposed cognitive mechanism's
representation, available operations, selection/control policy and
interpretation of background information. Do not use the DRR description
as a complete theorem prover, or use a formally valid conclusion to infer
that a participant applied a particular schema. Successful reading
performance and a mistake on a difficult reasoning task can both be
targets of one theory. The chapter's program-level descriptions and
empirical summaries therefore supplement, rather than replace, the
explicit evidence framework in \cref{sec:p19-evidence-methods}.

The P7 monotonicity treatment
(\cref{sec:p07-monotonicity-inference}) proves typed order-preserving
inference properties under specified semantic structures; it is another
use of ``natural logic,'' not the DRR cognitive architecture. Generic
ND, typed monotonicity, O'Brien's mental logic and a mental-model theory
must retain their different inputs, aims and evidence. The exact 1998
tables/program chapters and the outstanding named primary studies remain
necessary for a fully reconstructed inventory and empirical adjudication.
